\chapter{The Erd\H{o}s--Kac Theorem}

\begin{center}\itshape
Primes play a game of chance.\\
\upshape --- Mark Kac, \emph{Statistical Independence in Probability, Analysis and Number Theory} (1959)
\end{center}

How many prime factors does a large number have? The question sounds hopeless: $2^{100}$ has one, the product of the first twenty primes has twenty, and a
large prime has one. But ask instead about a \emph{typical} number, and the answer is astonishingly precise. In 1917 Hardy and Ramanujan proved that almost every
integer $n$ has about $\log\log n$ distinct prime factors. The function $\log\log n$ grows so slowly that this is a strong statement: for $n$ near $10^{9}$ it is $3.0$, and for $n$ with a
hundred digits it is only $5.4$. Their proof was long. In 1934 Pál Turán, then aged 24, found a proof in a single page: he computed a variance.

Then came one of the great stories of twentieth-century mathematics. In March 1939 Mark Kac gave a lecture at Princeton. He had been thinking about
independence, and he conjectured that the number of prime factors behaves like a sum of independent coin tosses, one for each prime $p$, coming up heads with probability
$\frac1p$. If so, the central limit theorem predicts a Gaussian law, with mean and variance $\log\log n$. Kac knew probability but lacked the number theory to prove it. Paul Erd\H{o}s was
in the audience. Kac recalled that before the lecture was over Erd\H{o}s had the proof, using the sieve. Their joint paper of 1940 founded probabilistic number theory.

\begin{think}
If $p\mid n$ were an independent event of probability $\frac1p$ for each prime $p\le n$, what would be the mean and the variance of the number of primes dividing $n$? (Mertens's second theorem,
Theorem~\ref{thm:mertens2}, gives the first.) Why should the variance be close to the mean?
\end{think}

This chapter proves the theorem of Erd\H{o}s and Kac in full. The proof is a dialogue between the integers and a probabilistic model in which divisibility by different primes is
independent, the model that Kubilius made precise in 1956. Turán's variance shows that the model has the right second moment. The model's higher moments are those of a
Gaussian, and the integers have the same moments as the model, because the small primes really are nearly independent. Moments then determine the law. Figure~\ref{fig:omega}
shows the data, and a surprise: at $10^7$ the bell curve is still far away.

\begin{figure}[ht]
\centering
\fig{ch09-omega}
\caption{The proportion of $n\le10^7$ with $\omega(n)=k$ (bars), against the Gaussian law of Erd\H{o}s and Kac with mean and variance $\log\log10^7=2.78$ (green), and the Poisson
law of mean $2.78$ shifted by one (orange dots). The data have mean $3.01$ but variance only $1.10$: the limit theorems hold, but $\log\log x$ grows so slowly that their limits are
never reached in practice.}
\label{fig:omega}
\end{figure}

\section{Turán's variance}

Write $\omega(n)$ for the number of distinct primes dividing $n$, and $\Omega(n)$ for the number counted with multiplicity, so $\omega(12)=2$ and $\Omega(12)=3$. Both are \emph{additive}:
$\omega(mn)=\omega(m)+\omega(n)$ when $\gcd(m,n)=1$. In fact $\omega$ is \emph{strongly additive}: $\omega(n)=\sum_{p\mid n}f(p)$ with $f(p)=1$. Chapter~1 found the mean (Exercise~\ref{ex:averageomega}):
$\sum_{n\le x}\omega(n)=x\log\log x+O(x)$. Turán's idea was to compute the variance too, and it works for every strongly additive function.

\begin{theorem}[the Turán--Kubilius inequality]\label{thm:tk}
There is an absolute constant $C$ such that, for every $x\ge3$ and every real function $f$ on the primes, the strongly additive function $f(n)=\sum_{p\mid n}f(p)$ satisfies
\[
\begin{gathered}\sum_{n\le x}\bigl(f(n)-A(x)\bigr)^2\le CxB(x)^2,\qq\text{where}\\A(x)=\sum_{p\le x}\frac{f(p)}p,\qq B(x)^2=\sum_{p\le x}\frac{f(p)^2}p.\end{gathered}
\]
\end{theorem}

In probabilistic language: for $n$ chosen uniformly in $[1,x]$, the variance of $f(n)$ is at most a constant times $B(x)^2$, which is (up to the factors $1-\frac1p$) the variance of
$\sum_pf(p)X_p$ for independent $X_p$ of mean $\frac1p$. Dependence can only cost a constant factor.

\begin{proof}
\emph{Small primes.} Suppose first that $f(p)=0$ for $p>\sqrt x$, and write $\sum_p$ for sums over $p\le\sqrt x$. Every $p\mid n$ with $n\le x$ is counted by $\lfloor\frac xp\rfloor$, so
\[
\begin{aligned}\sum_{n\le x}f(n)&=\sum_pf(p)\Bigl\lfloor\frac xp\Bigr\rfloor=xA+E_2,\eqbr\text{with }|E_2|\le\sum_p|f(p)|.\end{aligned}
\]
For the square, $f(n)^2=\sum_{p,q\mid n}f(p)f(q)$, and for $p\ne q$ both divide $n$ exactly when $pq$ does. Since $pq\le x$ for $p,q\le\sqrt x$,
\[
\begin{aligned}\sum_{n\le x}f(n)^2&=\sum_{p\ne q}f(p)f(q)\Bigl\lfloor\frac x{pq}\Bigr\rfloor+\sum_pf(p)^2\Bigl\lfloor\frac xp\Bigr\rfloor\\&\le x\Bigl(A^2-\sum_p\frac{f(p)^2}{p^2}\Bigr)+|E_1|+xB^2,\end{aligned}
\]
with $|E_1|\le\bigl(\sum_p|f(p)|\bigr)^2$. Expanding the square and using $\lfloor x\rfloor\le x$,
\[
\begin{aligned}\sum_{n\le x}(f(n)-A)^2&=\sum f(n)^2-2A\sum f(n)\eqbr+\lfloor x\rfloor A^2\\&\le xB^2+|E_1|+2|A||E_2|.\end{aligned}
\]
By Cauchy--Schwarz,
\[
\begin{aligned}\sum_p|f(p)|&\le\Bigl(\sum_p\frac{f(p)^2}p\Bigr)^{1/2}\Bigl(\sum_{p\le\sqrt x}p\Bigr)^{1/2}\eqbr\le B\sqrt{\sqrt x\,\pi(\sqrt x)},\end{aligned}
\]
so $|E_1|\le B^2\sqrt x\,\pi(\sqrt x)=O(\frac{xB^2}{\log x})$ by Chebyshev's bound
(Proposition~\ref{prop:chebyshev}), and $|E_2|\le Bx^{1/2}$. Also $|A|\le B\bigl(\sum_{p\le x}\frac1p\bigr)^{1/2}=O(B\sqrt{\log\log x})$, so $2|A||E_2|=O(B^2\sqrt x\log x)$. All three terms are $O(xB^2)$.

\emph{Large primes.} In general write $f=f_1+f_2$, where $f_1(p)=f(p)$ for $p\le\sqrt x$ and $f_2(p)=f(p)$ for $\sqrt x<p\le x$; then $A=A_1+A_2$, and $B_1,B_2\le B$. An integer $n\le x$ has at most
one prime factor $p>\sqrt x$, so $f_2(n)^2=f_2(p)^2$ or $0$, and $\sum_{n\le x}f_2(n)^2\le\sum_{\sqrt x<p\le x}f(p)^2\frac xp\le xB^2$. By Cauchy--Schwarz, $A_2^2\le B^2\sum_{\sqrt x<p\le x}\frac1p=O(B^2)$ by
Mertens's theorem. Hence $\sum_{n\le x}(f_2(n)-A_2)^2\le2\sum f_2(n)^2+2xA_2^2=O(xB^2)$, and $(a+b)^2\le2a^2+2b^2$ combines the two parts.
\end{proof}

\begin{theorem}[Hardy and Ramanujan 1917, Turán 1934]\label{thm:hr}
$\displaystyle\sum_{n\le x}\bigl(\omega(n)-\log\log x\bigr)^2=O(x\log\log x)$. Consequently, for every function $\psi(x)\to\infty$, the proportion of integers $n\le x$ with $|\omega(n)-\log\log x|>\psi(x)\sqrt{\log\log x}$ tends to $0$.
\end{theorem}
\begin{proof}
With $f(p)=1$, $A(x)=B(x)^2=\sum_{p\le x}\frac1p=\log\log x+O(1)$ by Theorem~\ref{thm:mertens2}. So $\sum_{n\le x}(\omega(n)-\log\log x)^2\le2\sum(\omega(n)-A)^2+2x(A-\log\log x)^2=O(x\log\log x)$. The second
statement is Chebyshev's inequality: the number of exceptions is at most $\frac{O(x\log\log x)}{\psi(x)^2\log\log x}$.
\end{proof}

One says that $\omega(n)$ has \emph{normal order} $\log\log n$: for every $\varepsilon>0$, the $n\le x$ with $|\omega(n)-\log\log n|>\varepsilon\log\log n$ have density $0$. (Replacing $\log\log x$ by $\log\log n$ costs nothing,
since $\log\log n=\log\log x+O(1)$ for $\sqrt x<n\le x$.) The same holds for $\Omega(n)$ (Exercise~\ref{ex:omegaOmega}), and this keeps a promise of Chapter~2.

\begin{corollary}\label{cor:taunormal}
For every $\varepsilon>0$, the integers $n$ with $(\log n)^{\log2-\varepsilon}\le\tau(n)\le(\log n)^{\log2+\varepsilon}$ have density $1$.
\end{corollary}
\begin{proof}
For every $n$, $2^{\omega(n)}\le\tau(n)\le2^{\Omega(n)}$, since $\tau(n)=\prod_{p^a\,\|\,n}(a+1)$ and $2\le a+1\le2^a$. For almost all $n$, both $\omega(n)$ and $\Omega(n)$ lie between $(1-\delta)\log\log n$ and $(1+\delta)\log\log n$,
and $2^{\log\log n}=(\log n)^{\log2}$.
\end{proof}

So a typical integer has about $(\log n)^{0.69}$ divisors, while the \emph{average} number of divisors is $\log n$ (Chapter~2). There is no contradiction: the average is dominated by rare integers with
very many divisors (Exercise~\ref{ex:taumean}). Means and typical values part ways whenever a quantity is the exponential of a sum, a lesson that probability teaches too.

\section{The model}

Kac's heuristic is a probability space on which divisibility by different primes is exactly independent. Fix $z\ge2$, and let $(X_p)_{p\le z}$ be independent random variables with
$P(X_p=1)=\frac1p$ and $P(X_p=0)=1-\frac1p$ (Volume~3 builds such finite families). The model of $\omega$ is $\sum_{p\le z}X_p$, and its centred version is
\[
\begin{gathered}S_z=\sum_{p\le z}\Bigl(X_p-\frac1p\Bigr),\qq\sigma_z^2=\operatorname{Var}S_z=\sum_{p\le z}\frac1p\Bigl(1-\frac1p\Bigr).\end{gathered}
\]
By Theorem~\ref{thm:mertens2}, $\sigma_z^2=\log\log z+O(1)\to\infty$. Since $S_z$ is a sum of many small independent pieces, it should be nearly Gaussian; we need this in the form of moments.
Let $m_k$ be the moments of the standard Gaussian law: $m_k=0$ for odd $k$ and $m_k=(k-1)!!=1\cdot3\cdots(k-1)$ for even $k$ (Volume~3).

\begin{lemma}\label{lem:modelmoments}
For every fixed $k\ge1$, $\mathbb E\bigl[(S_z/\sigma_z)^k\bigr]\to m_k$ as $z\to\infty$.
\end{lemma}
\begin{proof}
Write $Y_p=X_p-\frac1p$, so $\mathbb E Y_p=0$, $|Y_p|\le1$ and, for $j\ge2$, $|\mathbb E Y_p^j|\le\mathbb E Y_p^2=v_p$, where $v_p=\frac1p(1-\frac1p)\le\frac1p$. Expand
\[
\mathbb E[S_z^k]=\sum_{p_1,\dots,p_k\le z}\mathbb E[Y_{p_1}\cdots Y_{p_k}].
\]
Group the tuples $(p_1,\dots,p_k)$ according to the partition of $\{1,\dots,k\}$ into blocks of equal indices. If some block is a singleton $\{i\}$, then $Y_{p_i}$ is independent of the others and the term
vanishes. Consider a partition with blocks of sizes $j_1,\dots,j_b\ge2$, all at least $2$, so $b\le\frac k2$. By independence the term is $\prod_l\mathbb E[Y_{q_l}^{j_l}]$, for distinct primes $q_1,\dots,q_b$; its modulus is at
most $\prod_lv_{q_l}$, so the whole partition contributes at most $\bigl(\sum_{p\le z}v_p\bigr)^b=\sigma_z^{2b}$.

If $k$ is odd, every admissible partition has $b\le\frac{k-1}2$, and there are finitely many partitions (their number depends only on $k$), so $\mathbb E[S_z^k]=O_k(\sigma_z^{k-1})$, which is $o(\sigma_z^k)$.

If $k$ is even, the partitions with $b<\frac k2$ again contribute $O_k(\sigma_z^{k-2})$. The remaining ones are the pairings, all blocks of size $2$; there are $(k-1)!!$ of them (pair $1$ with one of the $k-1$ others,
and continue). Each pairing contributes $\sum v_{q_1}\cdots v_{q_{k/2}}$ over \emph{distinct} $q_1,\dots,q_{k/2}$. Without the distinctness condition the sum would be $\sigma_z^k$; the tuples with a repetition
contribute at most $\binom{k/2}2\bigl(\sum_pv_p^2\bigr)\sigma_z^{k-4}=O_k(\sigma_z^{k-4})$, since $\sum_pv_p^2\le\sum_p\frac1{p^2}<1$. So $\mathbb E[S_z^k]=(k-1)!!\,\sigma_z^k+O_k(\sigma_z^{k-2})$.
\end{proof}

Now the integers. For $n\ge1$ let $\omega_z(n)$ be the number of primes $p\le z$ dividing $n$, and $\mu_z=\sum_{p\le z}\frac1p$. Choose $n$ uniformly in $\{1,\dots,N\}$ and write $\mathbb E_N$ for the
average over $n\le N$.

\begin{lemma}\label{lem:compare}
For all $k\ge1$, $N\ge1$ and $z\ge2$,
\[
\Bigl|\mathbb E_N\bigl[(\omega_z(n)-\mu_z)^k\bigr]-\mathbb E[S_z^k]\Bigr|\le\frac{(2z)^k}N.
\]
\end{lemma}
\begin{proof}
Both sides are sums over $(p_1,\dots,p_k)$ with $p_i\le z$: on the left, $\omega_z(n)-\mu_z=\sum_{p\le z}(\mathbf 1_{p\mid n}-\frac1p)$. Fix such a tuple. Expanding the product,
\[
\prod_{i=1}^k\Bigl(\mathbf 1_{p_i\mid n}-\frac1{p_i}\Bigr)=\sum_{I\subset\{1,\dots,k\}}\ \prod_{i\notin I}\Bigl(-\frac1{p_i}\Bigr)\,\mathbf 1_{r_I\mid n},
\]
where $r_I$ is the product of the \emph{distinct} primes among the $p_i$, $i\in I$: the $p_i$ with $i\in I$ all divide $n$ exactly when $r_I$ does. The same expansion holds with $X_{p_i}$ in place of
$\mathbf 1_{p_i\mid n}$, and $\mathbb E\bigl[\prod_{i\in I}X_{p_i}\bigr]=\frac1{r_I}$ by independence and $X_p^2=X_p$. On the integers, $\mathbb E_N[\mathbf 1_{r\mid n}]=\frac1N\lfloor\frac Nr\rfloor=\frac1r-\frac\theta N$ with
$0\le\theta<1$. So the tuple contributes a difference of at most $\frac1N\sum_I\prod_{i\notin I}\frac1{p_i}=\frac1N\prod_i\bigl(1+\frac1{p_i}\bigr)\le\frac{2^k}N$. There are at most $z^k$ tuples.
\end{proof}

The lemma is where the independence of small primes becomes a theorem: as long as $z^k$ is small compared with $N$, the integers up to $N$ cannot tell the difference between
divisibility by the primes up to $z$ and independent coin tosses. This is Kubilius's model, and the reason it works is the reason the $2e^{-\gamma}$ of Chapter~1 appears: it works for the
small primes and fails for primes near $\sqrt N$.

\section{From moments to the bell curve}

The last ingredient is a principle of probability: if all the moments of $Y_N$ converge to those of a standard Gaussian, then $Y_N$ converges in law to it. Volume~3 states Lévy's
continuity theorem without proof; the special case we need is short, and we prove it from Volume~3's inversion formula. Let $\Phi$ be the distribution function of $N(0,1)$, and recall that
the characteristic function of $N(0,s^2)$ is $e^{-s^2t^2/2}$ (Volume~3, Proposition~6.7).

\begin{lemma}\label{lem:momentslaw}
Let $Y_N$ be real random variables with $\mathbb E[Y_N^k]\to m_k$ for every $k\ge1$. Then $P(Y_N\le a)\to\Phi(a)$ for every real $a$.
\end{lemma}
\begin{proof}
\emph{Characteristic functions converge.} For real $u$ and $K\ge1$, Taylor's formula with integral remainder gives $\bigl|e^{iu}-\sum_{k<K}\frac{(iu)^k}{k!}\bigr|\le\frac{|u|^K}{K!}$. Apply it with $u=tY_N$, take expectations,
and choose $K$ even, so that $\mathbb E|Y_N|^K=\mathbb E[Y_N^K]$:
\[
\Bigl|\varphi_{Y_N}(t)-\sum_{k<K}\frac{(it)^k}{k!}\mathbb E[Y_N^k]\Bigr|\le\frac{|t|^K}{K!}\mathbb E[Y_N^K].
\]
As $N\to\infty$ the sum tends to $\sum_{k<K}\frac{(it)^km_k}{k!}=\sum_{j<K/2}\frac{(-t^2/2)^j}{j!}$, because $\frac{m_{2j}}{(2j)!}=\frac1{2^jj!}$, and the right side tends to $\frac{|t|^K}{2^{K/2}(K/2)!}$. Both
$\bigl|\sum_{j<K/2}\frac{(-t^2/2)^j}{j!}-e^{-t^2/2}\bigr|$ and $\frac{|t|^K}{2^{K/2}(K/2)!}$ tend to $0$ as $K\to\infty$. So $\varphi_{Y_N}(t)\to e^{-t^2/2}$ for every $t$.

\emph{Smoothing.} Fix $\varepsilon>0$ and let $G$ be $N(0,1)$, independent of $Y_N$. The characteristic function of $Y_N+\varepsilon G$ is $\varphi_{Y_N}(t)e^{-\varepsilon^2t^2/2}$, which is integrable, so by the inversion formula
(Volume~3, Theorem~6.8(b)) $Y_N+\varepsilon G$ has the continuous density $q_N(y)=\frac1{2\pi}\int e^{-ity}\varphi_{Y_N}(t)e^{-\varepsilon^2t^2/2}\dd t$. By dominated convergence (the integrand is bounded by $e^{-\varepsilon^2t^2/2}$),
$q_N(y)\to q(y)$ for every $y$, where $q$ is the density of $N(0,1+\varepsilon^2)$. Then $\int|q_N-q|\to0$ (Scheffé's lemma: $\int|q-q_N|=2\int(q-q_N)^+$ since both integrate to $1$, and $0\le(q-q_N)^+\le q$, so dominated
convergence applies). Hence $P(Y_N+\varepsilon G\le b)\to\Phi\bigl(\frac b{\sqrt{1+\varepsilon^2}}\bigr)$ for every $b$.

\emph{Removing the smoothing.} For $\eta>0$, $\{Y_N\le a\}\subset\{Y_N+\varepsilon G\le a+\eta\}\cup\{\varepsilon G<-\eta\}$ and $\{Y_N+\varepsilon G\le a-\eta\}\subset\{Y_N\le a\}\cup\{\varepsilon G<-\eta\}$. So
\[
\begin{aligned}\liminf P(Y_N\le a)&\ge\Phi\Bigl(\frac{a-\eta}{\sqrt{1+\varepsilon^2}}\Bigr)-\Phi\Bigl(-\frac\eta\varepsilon\Bigr),\\\limsup P(Y_N\le a)&\le\Phi\Bigl(\frac{a+\eta}{\sqrt{1+\varepsilon^2}}\Bigr)+\Phi\Bigl(-\frac\eta\varepsilon\Bigr).\end{aligned}
\]
Let $\varepsilon\to0$ with $\eta$ fixed, then $\eta\to0$.
\end{proof}

\section{The theorem}

\begin{theorem}[Erd\H{o}s and Kac, 1940]\label{thm:erdoskac}
For every real $a$, the proportion of integers $n\le x$ such that
\[
\frac{\omega(n)-\log\log x}{\sqrt{\log\log x}}\le a
\]
tends to $\Phi(a)=\int_{-\infty}^a\frac{e^{-t^2/2}}{\sqrt{2\pi}}\dd t$ as $x\to\infty$.
\end{theorem}
\begin{proof}
Let $N=\lfloor x\rfloor$, $L=\log\log x$, and choose the cut-off $z=x^{1/\beta}$ with $\beta=L^{1/4}$, so that $\beta\to\infty$ but $\beta=o(\sqrt L)$.

\emph{The truncated function is Gaussian.} Let $Y_N(n)=\frac{\omega_z(n)-\mu_z}{\sigma_z}$ for $n$ uniform in $\{1,\dots,N\}$. By Lemma~\ref{lem:compare}, for each fixed $k$,
\[
\begin{aligned}\Bigl|\mathbb E_N[Y_N^k]-\mathbb E\Bigl[\Bigl(\frac{S_z}{\sigma_z}\Bigr)^k\Bigr]\Bigr|&\le\frac{(2z)^k}{N\sigma_z^k}\eqbr\le\frac{2^{k+1}x^{k/\beta}}x\to0,\end{aligned}
\]
since $\frac k\beta\to0$. With Lemma~\ref{lem:modelmoments}, $\mathbb E_N[Y_N^k]\to m_k$ for every $k$, and Lemma~\ref{lem:momentslaw} gives $P_N(Y_N\le b)\to\Phi(b)$ for every $b$.

\emph{Back to $\omega$.} For $n\le x$, the primes $p>z$ dividing $n$ number at most $\frac{\log x}{\log z}=\beta$, so $0\le\omega(n)-\omega_z(n)\le\beta$. By Theorem~\ref{thm:mertens2}, $\mu_z=\log\log z+M+o(1)=L-\log\beta+M+o(1)$, and
$\sigma_z^2=\mu_z+O(1)$, so $c_x=\frac{\sigma_z}{\sqrt L}\to1$. Hence
\[
\begin{gathered}\frac{\omega(n)-L}{\sqrt L}=c_xY_N(n)+\rho_x(n),\qq|\rho_x(n)|\le\eta_x=\frac{\beta+\log\beta+O(1)}{\sqrt L},\end{gathered}
\]
uniformly in $n\le x$, and $\eta_x\to0$. Since $c_x>0$, if $Y_N(n)\le\frac{a-\eta_x}{c_x}$ then $\frac{\omega(n)-L}{\sqrt L}\le a$, and if $\frac{\omega(n)-L}{\sqrt L}\le a$ then $Y_N(n)\le\frac{a+\eta_x}{c_x}$. Both bounds tend to $a$. So for every $\delta>0$
and $x$ large, the proportion in the theorem lies between $P_N(Y_N\le a-\delta)$ and $P_N(Y_N\le a+\delta)$, which tend to $\Phi(a-\delta)$ and $\Phi(a+\delta)$. Let $\delta\to0$.
\end{proof}

\begin{remark}[how fast, and beyond]
Rényi and Turán (1958) showed that the error in the theorem is $O\bigl(\frac1{\sqrt{\log\log x}}\bigr)$, and this is the true order. At $x=10^7$, $\frac1{\sqrt{\log\log x}}$ is still $0.6$: the proportion of $n\le10^7$ with $\omega(n)\le\log\log x$ is $0.32$, against the limit $0.5$, and Figure~\ref{fig:omega} looks nothing like a
bell curve. A sharper description is the Sathe--Selberg formula (1954): for fixed $k\ge1$, the number of $n\le x$ with $\omega(n)=k$ is
\[
\frac x{\log x}\frac{(\log\log x)^{k-1}}{(k-1)!}\Bigl(G\Bigl(\frac{k-1}{\log\log x}\Bigr)+o(1)\Bigr),
\]
for an explicit function $G$ with $G(0)=1$. Up to the factor $G$, this is the Poisson law of mean $\log\log x$, shifted by one, which is what the model predicts if one adds up rare independent
events; the case $G=1$ is Landau's theorem of 1900, met in Chapter~5. The Poisson law of large mean is close to Gaussian, which is the Erd\H{o}s--Kac theorem again. The same theorem holds for every
strongly additive $f$ with bounded values on the primes and $B(x)\to\infty$, with mean $A(x)$ and variance $B(x)^2$ (Kubilius, Shapiro): the proof above goes through with minor changes.
\end{remark}

\begin{beauty}
The theorem says that the primes dividing a random integer behave like independent coin tosses with probabilities $\frac12,\frac13,\frac15,\frac17,\dots$ The integers are as deterministic as anything
in mathematics, and the law of their prime factors is the law of chance. Kac called this ``statistical independence'', and it is the first theorem of this series in which number theory does not
only borrow tools from probability but obeys its laws.
\end{beauty}

\section{Exercises}

\begin{exercise}[$\star$]\label{ex:omegaOmega}
Show that $\sum_{n\le x}(\Omega(n)-\omega(n))=O(x)$. Deduce that $\Omega(n)$ has normal order $\log\log n$, and that Theorem~\ref{thm:erdoskac} holds with $\Omega$ in place of $\omega$.
\end{exercise}

\begin{exercise}[$\star$]\label{ex:taumean}
The average of $2^{\omega(n)}$ over $n\le x$ is about $\frac6{\pi^2}\log x$ (Exercise~\ref{ex:twoomega}), while a typical value is about $(\log x)^{0.69}$. Show that the $n\le x$ with $\omega(n)\le1.4\log\log x$
contribute $o(x\log x)$ to $\sum_{n\le x}2^{\omega(n)}$. So the sum is carried by a set of integers of density $0$.
\end{exercise}

\begin{exercise}[$\star$]\label{ex:skew}
Compute the third moment of $S_z$ and show that the skewness $\mathbb E[S_z^3]/\sigma_z^3$ is $\sim\frac1{\sqrt{\log\log z}}$. What does this say about the speed of convergence in Theorem~\ref{thm:erdoskac}?
\end{exercise}

\begin{exercise}[$\star\star$]\label{ex:largestprime}
Let $P(n)$ be the largest prime factor of $n$. Show that the proportion of $n\le x$ with $P(n)>\sqrt x$ tends to $\log2$. (The model of independent primes would answer $1-\prod_{\sqrt x<p\le x}(1-\frac1p)\to\frac12$. Why is the model wrong here?)
\end{exercise}

\begin{exercise}[$\star\star$, the signature exercise]\label{ex:multtable}
Let $D(N)$ be the number of distinct integers in the $N\times N$ multiplication table, that is, of the form $ab$ with $1\le a,b\le N$. Show that $D(N)=o(N^2)$. (Erd\H{o}s, 1955.)
\end{exercise}

\begin{exercise}[$\star\star\star$]\label{ex:omega1mod4}
Let $\omega_1(n)$ be the number of primes $p\equiv1\pmod4$ dividing $n$. Show that $\omega_1(n)$ has normal order $\frac12\log\log n$, and that $\frac{\omega_1(n)-\frac12\log\log x}{\sqrt{\frac12\log\log x}}$ is asymptotically $N(0,1)$
for $n\le x$.
\end{exercise}

\section{Hints}

\begin{hint}{ex:omegaOmega}
$\Omega(n)-\omega(n)=\sum_{p^a\|n}(a-1)\le\sum_{p^j\mid n,\,j\ge2}1$. Then Markov's inequality.
\end{hint}
\begin{hint}{ex:taumean}
Bound each term by its largest possible value, and use Theorem~\ref{thm:hr} for the density.
\end{hint}
\begin{hint}{ex:skew}
$\mathbb E[S_z^3]=\sum_p\mathbb E[Y_p^3]$ by the proof of Lemma~\ref{lem:modelmoments}.
\end{hint}
\begin{hint}{ex:largestprime}
An $n\le x$ has at most one prime factor $>\sqrt x$, so count the pairs $(p,m)$ with $n=pm$. For the model, which events are not independent?
\end{hint}
\begin{hint}{ex:multtable}
$\Omega(ab)=\Omega(a)+\Omega(b)$. Typical $a,b\le N$ have $\Omega(a)+\Omega(b)\approx2\log\log N$, but typical $m\le N^2$ have $\Omega(m)\approx\log\log N$.
\end{hint}
\begin{hint}{ex:omega1mod4}
Chapter~4 gives $\sum_{p\le x,\,p\equiv1\,(4)}\frac1p=\frac12\log\log x+O(1)$. Run the proofs of Theorems~\ref{thm:hr} and~\ref{thm:erdoskac} with only those primes.
\end{hint}

\section{Solutions}

\begin{solution}{ex:omegaOmega}
\idea Repeated prime factors are rare: $p^2\mid n$ has probability $\frac1{p^2}$, and $\sum\frac1{p^2}$ converges. So $\Omega$ and $\omega$ differ by a bounded amount on average, hence by a small amount almost always.
\writeup If $p^a\|n$ then $a-1$ is the number of $j\ge2$ with $p^j\mid n$. So $\sum_{n\le x}(\Omega(n)-\omega(n))\le\sum_{p^j\le x,\,j\ge2}\frac x{p^j}\le x\sum_p\frac1{p(p-1)}=O(x)$. By Markov's inequality, for $\psi(x)\to\infty$, the
$n\le x$ with $\Omega(n)-\omega(n)>\psi(x)$ number $O(\frac x{\psi(x)})=o(x)$. Taking $\psi=\varepsilon\log\log x$ gives the normal order. For Erd\H{o}s--Kac, take $\psi(x)=(\log\log x)^{1/4}$: outside a set of density $0$,
$\frac{\Omega(n)-\log\log x}{\sqrt{\log\log x}}$ and $\frac{\omega(n)-\log\log x}{\sqrt{\log\log x}}$ differ by at most $(\log\log x)^{-1/4}\to0$, and the end of the proof of Theorem~\ref{thm:erdoskac} applies.
\end{solution}

\begin{solution}{ex:taumean}
\idea A sum of exponentials is dominated by the largest exponents, even when they are rare. The typical $n$ has $2^{\omega(n)}\approx(\log x)^{0.69}$, far below the average, so the average must come from atypical $n$.
\writeup If $\omega(n)\le1.4\log\log x$, then $2^{\omega(n)}\le2^{1.4\log\log x}=(\log x)^{1.4\log2}=(\log x)^{0.9704}$. So these $n$ contribute at most $x(\log x)^{0.9704}=o(x\log x)$, while the whole sum is
$\frac6{\pi^2}x\log x+O(x)$. The rest of the sum comes from the $n$ with $\omega(n)>1.4\log\log x$, and by Theorem~\ref{thm:hr} (with $\psi(x)=0.4\sqrt{\log\log x}$) they have density $0$. The same argument with
$1.4$ replaced by any $\lambda<\frac1{\log2}=1.4427$ shows that the mean of $2^{\omega(n)}$ is made by integers with about $1.44\log\log x$ prime factors or more, far in the tail of the Erd\H{o}s--Kac law.
\end{solution}

\begin{solution}{ex:skew}
\idea Odd moments of the model come only from the blocks of size at least $3$; for the third moment, from triples of equal primes.
\writeup By independence and $\mathbb E Y_p=0$, $\mathbb E[S_z^3]=\sum_p\mathbb E[Y_p^3]$, and $\mathbb E[Y_p^3]=\frac1p(1-\frac1p)^3-(1-\frac1p)\frac1{p^3}=\frac1p(1-\frac1p)(1-\frac2p)$. So $\mathbb E[S_z^3]=\log\log z+O(1)$, like $\sigma_z^2$,
and the skewness is $\frac{\log\log z+O(1)}{(\log\log z+O(1))^{3/2}}\sim\frac1{\sqrt{\log\log z}}$. For a sum of independent variables, the Berry--Esseen theorem (stated in Volume~3) bounds the distance to the Gaussian law by a constant times $\sum_p\mathbb E|Y_p|^3/\sigma_z^3$, which is of the same order $\frac1{\sqrt{\log\log z}}$; Rényi and Turán proved the same rate for $\omega$ itself. With all primes up to $10^7$
the model has skewness $0.49$, and $\log\log z$ would have to exceed $10^4$, so $z>e^{e^{10^4}}$, to bring it below $0.01$.
\end{solution}

\begin{solution}{ex:largestprime}
\idea Since at most one prime factor exceeds $\sqrt x$, the events ``$p\mid n$'' for $p>\sqrt x$ are disjoint, not independent. So their probabilities add up, to $\sum_{\sqrt x<p\le x}\frac1p\to\log2$, instead of combining
like independent events.
\writeup If $n\le x$ has a prime factor $p>\sqrt x$, it has exactly one, since the product of two would exceed $x$. So the number of such $n$ is $\sum_{\sqrt x<p\le x}\lfloor\frac xp\rfloor=x\sum_{\sqrt x<p\le x}\frac1p+O(\pi(x))$.
By Theorem~\ref{thm:mertens2}, $\sum_{\sqrt x<p\le x}\frac1p=\log\log x-\log\log\sqrt x+O(\frac1{\log x})=\log2+O(\frac1{\log x})$, and $\pi(x)=O(\frac x{\log x})$. The proportion tends to $\log2=0.693$. The model treats ``$p\mid n$'' and ``$q\mid n$'' as independent for $p,q>\sqrt x$, but they are
incompatible; its answer is $1-\exp\bigl(-\sum\frac1p+O(\frac1{\sqrt x})\bigr)\to1-\frac12=\frac12$. This is the same failure of the model for large primes as the $2e^{-\gamma}$ of Chapter~1, and the reason Lemma~\ref{lem:compare} needs $z$ small.
\end{solution}

\begin{solution}{ex:multtable}
\idea Products know how many prime factors their factors had. A typical product $ab$ has about $2\log\log N$ prime factors, but a typical number up to $N^2$ has only about $\log\log N^2\approx\log\log N$. So the typical products are
atypical numbers, and atypical numbers are few.
\writeup Fix $\varepsilon\in\left]0,\frac14\right[$ and let $L=\log\log N$. By Exercise~\ref{ex:omegaOmega}, all but $o(N)$ integers $a\le N$ have $|\Omega(a)-L|\le\varepsilon L$. So all but $o(N^2)$ pairs $(a,b)$ have
$\Omega(ab)=\Omega(a)+\Omega(b)\ge(2-2\varepsilon)L$. On the other hand, $\log\log N^2=L+\log2$, so by the normal order of $\Omega$ (with $x=N^2$), all but $o(N^2)$ integers $m\le N^2$ have $\Omega(m)\le(1+\varepsilon)L+\log2<(2-2\varepsilon)L$ for large $N$.
Every product $ab$ is either the product of an exceptional pair, or an integer $m\le N^2$ with $\Omega(m)\ge(2-2\varepsilon)L$, an exceptional integer. So $D(N)\le o(N^2)+o(N^2)$.
Numerically, one finds $D(N)/N^2=0.29$, $0.25$, $0.24$ for $N=100$, $1000$, $2000$: the decrease is very slow. Ford (2008) found the true order, $D(N)\asymp\frac{N^2}{(\log N)^c(\log\log N)^{3/2}}$ with $c=1-\frac{1+\log\log2}{\log2}=0.086\ldots$
\end{solution}

\begin{solution}{ex:omega1mod4}
\idea Nothing in the proofs used that \emph{all} primes are counted: only the sum of $\frac1p$ over the primes used, which here is $\frac12\log\log x+O(1)$.
\writeup Apply Theorem~\ref{thm:tk} to $f(p)=1$ for $p\equiv1\pmod4$ and $f(p)=0$ otherwise: $A(x)=B(x)^2=\sum_{p\le x,\,p\equiv1\,(4)}\frac1p=\frac12\log\log x+O(1)$ by Corollary~\ref{cor:recipprog}. As in Theorem~\ref{thm:hr},
$\sum_{n\le x}(\omega_1(n)-\frac12\log\log x)^2=O(x\log\log x)$, which gives the normal order. For the Gaussian law, repeat Lemmas~\ref{lem:modelmoments} and~\ref{lem:compare} with the primes $p\le z$, $p\equiv1\pmod4$ only:
the model has variance $\sigma_z^2=\sum_{p\le z,\,p\equiv1}\frac1p(1-\frac1p)=\frac12\log\log z+O(1)\to\infty$, which is all that Lemma~\ref{lem:modelmoments} used, and Lemma~\ref{lem:compare} holds verbatim. The end of the proof of
Theorem~\ref{thm:erdoskac} goes through with $L$ replaced by $\frac12\log\log x$: the truncation error is still at most $\beta$, and $\mu_z=\frac12\log\log x-\frac12\log\beta+O(1)$.
\end{solution}

\section*{Sources and further reading}

\begin{itemize}
\item G.~H.~Hardy and S.~Ramanujan, \emph{The normal number of prime factors of a number $n$}, Quarterly Journal of Mathematics 48 (1917), 76--92.
\item P.~Turán, \emph{On a theorem of Hardy and Ramanujan}, Journal of the London Mathematical Society 9 (1934), 274--276.
\item P.~Erd\H{o}s and M.~Kac, \emph{The Gaussian law of errors in the theory of additive number theoretic functions}, American Journal of Mathematics 62 (1940), 738--742.
\item M.~Kac, \emph{Statistical Independence in Probability, Analysis and Number Theory}, Carus Mathematical Monographs 12, MAA, 1959. Short, and a joy to read.
\item M.~Kac, \emph{Enigmas of Chance: an Autobiography}, Harper and Row, 1985. The story of the 1939 lecture.
\item J.~Kubilius, \emph{Probabilistic Methods in the Theory of Numbers}, AMS Translations of Mathematical Monographs 11, 1964 (Russian original 1959).
\item A.~Rényi and P.~Turán, \emph{On a theorem of Erd\H{o}s--Kac}, Acta Arithmetica 4 (1958), 71--84.
\item A.~Selberg, \emph{Note on a paper by L.~G.~Sathe}, Journal of the Indian Mathematical Society 18 (1954), 83--87.
\item P.~Billingsley, \emph{On the central limit theorem for the prime divisor function}, American Mathematical Monthly 76 (1969), 132--139. The method of moments, as used here.
\item A.~Granville and K.~Soundararajan, \emph{Sieving and the Erd\H{o}s--Kac theorem}, in \emph{Equidistribution in Number Theory, an Introduction}, Springer, 2007, 15--27.
\item K.~Ford, \emph{The distribution of integers with a divisor in a given interval}, Annals of Mathematics 168 (2008), 367--433. The multiplication table problem.
\item G.~Tenenbaum, \emph{Introduction to Analytic and Probabilistic Number Theory}, 3rd ed., AMS Graduate Studies in Mathematics 163, 2015, Part~III.
\end{itemize}
